\paragraph{最初的结果。}\label{sub:first-consequences}
\cref{thm:hurwitz-pi-twist-isomorphic}的最直接的一个应用，是把\trHurwitz{}移动与\trMutation{}联系起来。取\trVanishingPath{}的一个\trDistinguishedBasis{}$\boldsymbol{\gamma}=(\gamma_1,\ldots,\gamma_m)$，并考虑与之相伴的 (B) 型\trEquivariant\trBrane{}的集合 $\widetilde{\boldsymbol{\Delta}}^{\#}_{\boldsymbol{\gamma}}$——其中诸\trGrading{}的选取是任意的——以及相应的\trVanishingCycle{}集合 $\boldsymbol{V}^{\#}_{\boldsymbol{\gamma}}$，后者带有诱导的\trBraneStructure。\cref{thm:fibration-exceptional-collection}已经告诉我们，$\widetilde{\boldsymbol{\Delta}}^{\#}_{\boldsymbol{\gamma}}$ 是一个\trExceptionalCollection。而\cref{thm:hurwitz-pi-twist-isomorphic}则说：如果把 \eqref{eq:hurwitz-move-spheres} 所定义的初等\trHurwitz{}移动施加到我们的基上，那么所得结果就是该\trExceptionalCollection{}的一个初等\trMutation{}（在 \eqref{eq:elementary-mutation} 的意义下）。我们在选取相关记号时其实已经暗中认可了这种对应，所以可以写成
\begin{equation}\label{eq:hurwitz-mutation-iso}
\widetilde{\boldsymbol{\Delta}}^{\#}_{\mathscr{L}_k\boldsymbol{\gamma}}\cong\mathscr{L}_k\widetilde{\boldsymbol{\Delta}}^{\#}_{\boldsymbol{\gamma}}.
\end{equation}
严格说来，要使 \eqref{eq:hurwitz-mutation-iso} 成为 $D\scr{F}(\pi)$ 中\trExceptionalCollection{}的一个真正的\trIsomorphism，诸\trGrading{}必须像\cref{thm:hurwitz-pi-twist-isomorphic}中那样具体地选取（否则，它将只是一个相差\trTranslation{}的\trIsomorphism，不过这对我们的目的来说已经完全够用了）。当然，只需改变所涉集合的记号，同样的结果对逆 $\mathscr{L}_k^{-1}$ 也成立，并且经\trComposition{}后对任何\trMutation{}都成立。换言之，\trBraidGroup{} $\mathrm{Br}_m$ 在\trExceptionalCollection{}上定义的那个代数作用（相差\trTranslation{}与\trIsomorphism；见\cref{sub:mutation}）可以与它在\trVanishingPath{}的\trDistinguishedBasis{}上的几何作用（见\cref{sub:distinguished-bases}）相等同。
取定某个 $\boldsymbol{\gamma}$。暂时用 $\scr{T}$ 表示 $Tw\scr{F}(\pi)$ 中由诸\trObject{}$\widetilde{\Delta}^{\#}_{\gamma_1},\ldots,\widetilde{\Delta}^{\#}_{\gamma_m}$ 生成的\trTriangulated{}子\trCategory。不言自明地，$\widetilde{\boldsymbol{\Delta}}^{\#}_{\boldsymbol{\gamma}}$ 在该子\trCategory{}中是一个\trFull{} \trExceptionalCollection，而且这一点在一切\trMutation{}之下继续保持。利用 \eqref{eq:hurwitz-mutation-iso}，这意味着：如果对基施加一次\trHurwitz{}移动并考虑与之相伴的集合，那么它的所有\trObject{}仍落在 $\scr{T}$ 之中，并且生成该\trCategory。由于\trHurwitz{}移动在诸例外基上的作用是可递的，由此可知 $\scr{T}$ 与 $\boldsymbol{\gamma}$ 无关。事实上，\cref{thm:type-u-in-generated-subcategory}告诉我们，任何 (U) 型\trEquivariant\trBrane{}都落在 $\scr{T}$ 中，于是按 $\scr{F}(\pi)$ 的定义必有 $\scr{T}=Tw\scr{F}(\pi)$。另一方面，\cref{thm:directed-derived-equivalence}说，$\scr{T}$\trQuasiequivalent{}于与 $\widetilde{\boldsymbol{\Delta}}^{\#}_{\boldsymbol{\gamma}}$ 相伴的\trDirected{} \trAinfty-子\trCategory{}之上的\trTwistedComplex\trCategory，而后者又由\cref{thm:fibration-exceptional-collection}\trQuasiisomorphic{}于 $\scr{F}^{\rightarrow}(\boldsymbol{V}^{\#}_{\boldsymbol{\gamma}})$。于是：
\begin{theorem}\label{thm:tw-directed-quasiequiv}
$Tw\scr{F}^{\rightarrow}(\boldsymbol{V}^{\#}_{\boldsymbol{\gamma}})$ 与 $Tw\scr{F}(\pi)$\trQuasiequivalent{}。
\end{theorem}
\noqed
作为一个直接的推论，我们得到\cref{thm:twisted-directed-fibration-invariant}的另一个证明（而且事实上\cref{thm:twisted-directed-hurwitz-invariant}也是如此，因为任何一组\trLagrange\trSphere{}都可以实现为一个适当构造的\trLefschetzFibration{}的\trVanishingCycle{}基）。这实际上是对原来那个证明的一个改进：通过把 $Tw\scr{F}^{\rightarrow}(\boldsymbol{V}^{\#}_{\boldsymbol{\gamma}})$ 等同于 $Tw\scr{F}(\pi)$，便证明了它与 $\boldsymbol{\gamma}$ 无关，而后者显然不涉及基的任何选取。作为一个具体的收益，我们可以对不同 $\boldsymbol{\gamma}$ 的诸\trCategory{}$Tw\scr{F}^{\rightarrow}(\boldsymbol{V}^{\#}_{\boldsymbol{\gamma}})$ 之间构造出\trCanonical{}的\trQuasiequivalence——至少在相应\trFunctor{}\trCategory{}$H^{0}(fun)$ 中相差一个\trIsomorphism{}的意义下是\trCanonical{}的。
为此所付出的唯一代价是技术性假设 $\mathrm{char}(\KK)\neq 2$，它是使用\trDoubleBranchedCover{}技巧而强加给我们的。其次，注意\cref{thm:full-embedding-functor}与\cref{thm:tw-preserves-cohomologically-ff}一起给出一个\trCohomology{}\trFullyFaithful{}的\trFunctor{}$Tw\scr{F}(E)\to Tw\scr{F}(\pi)$。鉴于\cref{thm:tw-directed-quasiequiv}，这就给出：
\begin{corollary}\label{thm:directed-fukaya-ff-embedding}
对任何 $\boldsymbol{\gamma}$，存在 $\scr{F}(E)$ 到 $Tw\scr{F}^{\rightarrow}(\boldsymbol{V}^{\#}_{\boldsymbol{\gamma}})$ 中的一个\trCohomology{}\trFullyFaithful{}的\trEmbedding。
\end{corollary}
\noqed
要点在于：既然\trDirected{} \trFukayaCategory{}是在 $\scr{F}(M)$ 之内构造的，那么从抽象的角度说，\trTotalSpace{} $E$ 中的\trFloerCohomology{}计算就可以完全化归为在\trFiber{}中进行的计算。下面我们尝试使这些考虑具有一个更具体的形状。
